\documentclass[12pt,a4paper]{article}
\usepackage[margin=25mm, headheight=15pt, headsep=10pt]{geometry}
\usepackage{fontspec}
\usepackage[table]{xcolor} % Note: table option is required for rowcolors
\usepackage{titlesec}
\usepackage{tocloft}
\usepackage{fancyhdr}
\usepackage{booktabs}
\usepackage{tcolorbox}
\tcbuselibrary{skins, breakable}
\usepackage[colorlinks=true, linkcolor=emerald, urlcolor=emerald, bookmarks=true]{hyperref}

\definecolor{emerald}{RGB}{0,102,51}
\definecolor{darkred}{RGB}{139,0,0}
\definecolor{lightgray}{RGB}{245,245,245}

\usepackage[nil]{babel}
\babelprovide[import, main]{bengali}
\babelprovide[import]{arabic}
\babelfont{rm}[Renderer=HarfBuzz,Scale=1.0]{Noto Serif Bengali}
\babelfont{sf}[Renderer=HarfBuzz]{Noto Sans Bengali}
\babelfont{tt}[Renderer=HarfBuzz]{Noto Sans Bengali}
\babelfont[arabic]{rm}[Renderer=HarfBuzz,Scale=1.1]{Noto Naskh Arabic}

% Section styling
\titleformat{\section}{\Large\bfseries\color{emerald}}{\thesection.}{0.5em}{}
\titleformat{\subsection}{\large\bfseries\color{emerald!85!black}}{\thesubsection.}{0.5em}{}
\titleformat{\subsubsection}{\normalsize\bfseries\color{emerald!70!black}}{\thesubsubsection.}{0.5em}{}
\titlespacing*{\section}{0pt}{18pt}{8pt}

% Fancy Header/Footer setup
\pagestyle{fancy}
\fancyhf{} % clear all header and footer fields
\fancyhead[L]{\color{emerald}\small\textbf{\nouppercase{\leftmark}}}
\fancyfoot[L]{\scriptsize\color{black!50}{\lat{AI}}-সংকলিত, আলেম-যাচাইকৃত নয়}
\fancyfoot[C]{\color{emerald!70!black}\thepage}
\renewcommand{\headrulewidth}{0.4pt}
\renewcommand{\headrule}{\hbox to\headwidth{\color{emerald}\leaders\hrule height \headrulewidth\hfill}}
\renewcommand{\footrulewidth}{0pt}

% Custom boxes
% 1. Ayat/Hadith Box
\newtcolorbox{ayatbox}[1][]{
    enhanced, breakable, 
    colback=emerald!5, 
    colframe=emerald, 
    boxrule=0.5pt, 
    arc=3pt, 
    left=10pt, right=10pt, top=10pt, bottom=10pt,
    #1
}

% 2. Quote/Translation Box
\newtcolorbox{quotebox}[1][]{
    enhanced, breakable,
    frame hidden,
    colback=lightgray,
    borderline west={3pt}{0pt}{emerald},
    left=15pt, right=10pt, top=10pt, bottom=10pt,
    sharp corners,
    #1
}

% 3. AI-generated notice box
\newtcolorbox{warnbox}[1][]{
    enhanced, breakable,
    colback=red!4,
    colframe=darkred,
    boxrule=0.8pt,
    arc=3pt,
    left=10pt, right=10pt, top=10pt, bottom=10pt,
    #1
}

% Commands
\newcommand{\arab}[1]{\foreignlanguage{arabic}{#1}}
\newcommand{\kib}[1]{\textcolor{emerald}{\textbf{#1}}}

% Latin (English) fallback: Noto Serif Bengali has NO Latin glyphs
\newfontfamily\latfont[Renderer=HarfBuzz]{Noto Serif}
\newcommand{\lat}[1]{{\latfont #1}}

\setlength{\parskip}{5pt}
\setlength{\parindent}{0pt}

% Fix for missing bullet in Noto Serif Bengali
\renewcommand{\labelitemi}{$\bullet$}



\begin{document}
\begin{center}{\Large\bfseries\color{emerald} \lat{01}-নবীজির-খুশু-নামাজ}\end{center}
\vspace{1em}
\section*{০১ — নবীজি \arab{ﷺ}: রাতের নামাজে পা ফুলে যাওয়া}

\begin{warnbox}
 \textbf{গুরুত্বপূর্ণ নোটিশ:} এই কন্টেন্টটি \lat{AI} দিয়ে প্রস্তুত। কেবল সহীহ/হাসান হাদিস রাখা হয়েছে।
\end{warnbox}

\medskip\hrule\medskip

\subsection*{১. সূত্র কার্ড}

\begin{table}[h]\centering\small
\begin{tabular}{ll}
বিষয় & বিবরণ \\
\hline
\textbf{বর্ণনাকারী} & আয়েশা (রাঃ) \\
\textbf{গ্রন্থ} & সুনানে তিরমিযী ২৪১৮, সুনানে আবু দাউদ ১৩১২ (সহীহ — আলবানি) \\
\textbf{মূল বিষয়} & নবীজি \arab{ﷺ}-এর রাতের নামাজের তীব্রতা ও কৃতজ্ঞতা \\
\hline
\end{tabular}
\end{table}

\medskip\hrule\medskip

\subsection*{২. পটভূমি}

মদীনায় রাতের নামাজ (তাহাজ্জুদ) — নবীজি \arab{ﷺ} প্রতিদিন রাতে উঠে নামাজ পড়তেন। আয়েশা (রাঃ) তাঁকে রাতের নামাজে দেখেছেন — পায়ের সোল ফুলে গিয়েছিল, কিন্তু তিনি থামেননি।

\medskip\hrule\medskip

\subsection*{৩. পূর্ণ ঘটনার বর্ণনা}

আয়েশা (রাঃ) বলেন:

\begin{quote}
\textbf{\lat{“}একবার আমি নবীজি (সা.)-কে বললাম: "হে আল্লাহর রাসূল! আপনি এত বেশি ইবাদত করছেন যে আপনার পায়ের সোল ফুলে গেছে। আপনি কি এমন করছেন যে আল্লাহ আপনাকে ক্ষমা করবেনই?"} \\
\textbf{নবীজি (সা.) বললেন: "আমি কি একজন কৃতজ্ঞ বান্দা হতে পারি না?"\lat{”}}
\end{quote}

\medskip\hrule\medskip

\subsection*{৪. শিক্ষা}

\subsubsection*{\textbf{"কৃতজ্ঞ বান্দা" — নবীজি \arab{ﷺ}-এর মানদণ্ড}}
নবীজি \arab{ﷺ}-এর জন্য ইবাদত শুধু \textbf{নফর} (কর্তব্য) নয় — বরং \textbf{\lat{shukr} (কৃতজ্ঞতা)}। তিনি আল্লাহর অনুগ্রহ স্মরণ করে এত পড়তেন যে পা ফুলে যেত — এটি \textbf{\lat{love} (ভালোবাসা)} ও \textbf{\lat{gratitude} (কৃতজ্ঞতা)}-র চূড়ান্ত রূপ।

\subsubsection*{\textbf{"আমি কি একজন কৃতজ্ঞ বান্দা হতে পারি না?"}}
এই প্রশ্ন \textbf{\lat{humility} (বিনয়)}-এর উদাহরণ। নবীজি \arab{ﷺ} নিজেকে \textbf{"কৃতজ্ঞ বান্দা"} বলতে চাইছিলেন — তাঁর জন্য ইবাদত ছিল \textbf{\lat{reward} (পুরস্কার)} নয়, বরং \textbf{\lat{gift} (দান)}-এর প্রতিদান।

\medskip\hrule\medskip

\subsection*{৫. সংশ্লিষ্ট হাদিস}

\begin{table}[h]\centering\small
\begin{tabular}{ll}
হাদিস & গ্রন্থ \\
\hline
\textbf{\lat{“}আমি কি একজন কৃতজ্ঞ বান্দা হতে পারি না?\lat{”}} & সুনানে তিরমিযী ২৪১৮ \\
\textbf{\lat{“}নবীজি (সা.) রাতে জেগে থাকতেন — তাঁর পায়ের সোল ফুলে যেত।\lat{”}} & সুনানে আবু দাউদ ১৩১২ \\
\hline
\end{tabular}
\end{table}

\medskip\hrule\medskip

\subsection*{৬. প্রয়োগ}

নামাজে যখন \textbf{\lat{tired} (ক্লান্ত)} বোধ করেন — মনে করুন: \textbf{"আমি কি একজন কৃতজ্ঞ বান্দা হতে পারি না?"} এটি আপনাকে \textbf{\lat{motivated} (অনুপ্রাণিত)} করবে।

\medskip\hrule\medskip

\textit{শেষ আপডেট: আগস্ট ২০২৬}

\end{document}
